\begin{helper}[Base of a nontrivial ray]
Let $F$ be a nontrivial front on $X$, let $n\in X$, and suppose
$\mathsf{Ray}(F,n)\ne\{\emptyset\}$.  Then

\[
\mathsf{FrontBase}(\mathsf{Ray}(F,n))=\mathsf{TailSet}(X,n).
\]
\end{helper}
\begin{proof}
For the inclusion from left to right, let
$k\in\mathsf{FrontBase}(\mathsf{Ray}(F,n))$.  By
\noderef{FrontBase}, choose $s\in\mathsf{Ray}(F,n)$ with $k\in s$.
The definition of \noderef{Ray} gives $n<k$ and
$\{n\}\cup s\in F$.  Thus $k\in\mathsf{FrontBase}(F)$, so
$k\in X$ by \noderef{FrontNontrivialBase}; consequently
$k\in\mathsf{TailSet}(X,n)$ by \noderef{TailSet}.

Conversely, fix $k\in\mathsf{TailSet}(X,n)$ and put
$Y=\{k\}\cup\mathsf{TailSet}(X,k)$.  The tail
$\mathsf{TailSet}(X,k)$ is infinite: it is obtained from the infinite
base $X$ by removing the finite initial segment $\{0,\ldots,k\}$.  Hence
$Y$ is infinite.  Also $Y\subseteq\mathsf{TailSet}(X,n)$: its first
element is $k$, and every element of its second part is greater than $k$.
By \noderef{FrontRayDense}, choose $s\in\mathsf{Ray}(F,n)$ and a cutoff
$m\in Y$ such that

\[
s=\{j\in Y\mid j<m\}.
\]

The set $s$ is nonempty.  Otherwise $\emptyset\in\mathsf{Ray}(F,n)$;
then every $t\in\mathsf{Ray}(F,n)$ is equal to $\emptyset$, because
$\emptyset$ is an initial segment of $t$ by \noderef{InitialSegment}, and
\noderef{FrontRayPrefixFree} applies.  This would make the ray trivial,
contrary to the hypothesis.

Choose $q\in s$.  Since $q\in Y$, either $q=k$ or $k<q$, so $k\le q<m$.
Therefore $k<m$, and the displayed cutoff equation gives $k\in s$.
The definition of \noderef{FrontBase} now yields
$k\in\mathsf{FrontBase}(\mathsf{Ray}(F,n))$, proving the reverse
inclusion.
\end{proof}
